\documentclass[10pt]{article}
\usepackage[margin=20mm]{geometry}
\usepackage{amsmath,amssymb,mathtools,fontspec,unicode-math,xurl,hyperref}
\IfFontExistsTF{DejaVu Serif}{\setmainfont{DejaVu Serif}}{\setmainfont{Latin Modern Roman}}
\setmathfont{latinmodern-math.otf}
\setlength{\parindent}{0pt}\setlength{\parskip}{6pt}\sloppy
\title{Initial scientific review of
Global solutions of the two-dimensional Riemann problem with four-shock interactions for the Euler equations for potential flow}
\author{}
\date{Frozen batch M25-0249; initial review, decision pending}
\begin{document}
\maketitle
Abstract: The initial assessment is Major. The coefficient \(O_3\) in the sonic-coordinate equation is algebraically inconsistent with the original potential-flow equation. An exact, strictly elliptic local solution gives residual \(1/192\) in the printed transformed equation. The same coefficient prevents the cancellation of lower-order terms claimed in constructing the iteration operator. This affects a core proof dependency, even though a concrete correction is available and the leading sonic degeneracy appears unchanged. Two further Major derivation errors affect the radial comparison used in the subsonic-corner hodograph argument and the twice-differentiated identity used for shock joining in Theorem 2.2(i). Their corrected arguments restore the specified local steps. Four other groups of bounded errors concern mapping estimates, vacuum, the normal endpoint and Bernoulli algebra. The global existence, regularity and convexity theorems themselves have not been disproved. Scientific decision and adjudication remain pending.\par
Article and binding. Gui-Qiang G. Chen, Alex Cliffe, Feimin Huang, Song Liu and Qin Wang. Frozen journal identifier: 10.4171/JEMS/1622.\par
Version examined. The acquired file is arXiv:2305.15224v1, submitted 24 May 2023, with PDF date 25 May 2023 and 75 pages. Page numbers below are both printed and 1-based PDF page numbers. The frozen title uses of potential flow and the frozen author field is Alexander J. Cliffe; the actual PDF uses for potential flow and Alex Cliffe. These fields are retained separately. The publisher records online publication on 14 April 2025 and requires a subscription. No equivalence to the journal version has been established. All findings below concern the acquired v1.\par
\section*{Scope, hypotheses and dependency chain}
The primary reviewer fully read all 75 source-text pages and actually inspected 55 rendered source pages, including every page used below as evidence. Their text ranges, absence of independent image inspection, and developer-provided model self-reports are recorded in the supporting files. The primary reviewer independently evaluated the calculations and arguments reported here. This is initial scientific review, not independent acceptance or human review.\par
The model is compressible, irrotational, isentropic potential flow. It is not the full Euler system with arbitrary vorticity. In self-similar coordinates, the pseudo-potential \(\varphi\) satisfies \[\operatorname{div}(\rho\mathrm D\varphi)+2\rho=0,\qquad
 \rho=\left(1+(\gamma-1)
       \left(B-\frac{|\mathrm D\varphi|^2}{2}-\varphi\right)\right)^{1/(\gamma-1)},\] with the exponential interpretation for \(\gamma=1\). The physical self-similar potential is \(\phi=\varphi+|\boldsymbol\xi|^2/2\). The normalization is \(\rho_2=\rho_4=1\), \(u_2=v_1=0\), \(B=v_2^2/2\) and \(v_{\min}<v_2<0\). The two incident angle pairs are related by reflection symmetry, leaving two independent angles, not four arbitrary independent incoming shocks.\par
Compatibility fixes the incident densities using \[\ell(\rho,1)=\sqrt{\frac{2(\rho-1)h(\rho)}{\rho+1}},
 \qquad
 \ell(\rho_1,1)\cos\theta_1
 =\ell(\rho_3,1)\cos\theta_2=-v_2.\] The vacuum critical angle, sonic angle and detachment angle have distinct roles. The theorem’s parameter region \(\Theta=[0,\theta^d)^2\setminus\{(0,0)\}\) includes unilateral zero angles and mixed reflection types. Normal reflection is an explicitly supplied limiting configuration and a base point for the degree argument. Sonic arcs collapse in the appropriate subsonic cases. These boundary cases were retained in the review.\par
Theorem 2.1 claims existence and convergence to normal reflection, not unrestricted uniqueness. Theorem 2.2 separates \(C^{1,1}\) regularity near noncollapsed sonic arcs from \(C^{1,\alpha}\) corner regularity and specifies the second-derivative jump and endpoint failure of a Hessian limit. Theorem 2.3 concerns strict convexity and uniform convexity on compact subarcs of the relative interior; the flat normal shock is outside \(\Theta\).\par
The proof proceeds through directional monotonicity and a priori estimates, moving-domain flattening, an iteration map on weighted spaces, a fixed-point index at normal reflection, and a continuation argument. At a fixed point, the shock graphs must agree and the modified equation must reduce to the original one. F1 concerns precisely this latter identification. The review also checked the normal-reflection formulas, the positivity range of the incident densities, the logic of the directional-derivative contradiction in Lemma 3.1, and the scaling in the four sonic/corner regimes. General Schauder estimates and the full external existence and convexity results cited as [1], [2], [10], [12] and [48] were not independently reproduced.\par
\section*{Confirmed issues}
\subsection*{F1. An incorrect sonic coefficient changes the iteration equation (Major)}
Location. Equations (3.24)–(3.26), pp. 25–26, especially \(O_3\) on p. 26; the modified coefficients and assertion that Cartesian lower-order terms vanish, p. 46; the fixed-point identification in Proposition 4.23, Step 3, p. 57.\par
Write \(r=c-x\), \(c=c_5\), \(p=\psi_x\), \(q=\psi_y\), \(z=\psi\), where \(\psi=\varphi-\varphi_5\) in the sonic polar coordinates. The actual squared sound speed is \[a=c^2-(\gamma-1)
 \left(z+rp+\frac{p^2}{2}+\frac{q^2}{2r^2}\right).\] The polar components of \(\mathrm D\varphi\) are \((-r-p,q/r)\). Direct substitution into the nondivergence form of the original equation, followed by division by \(c\), gives \[\begin{aligned}
 &\frac{a-(r+p)^2}{c}\psi_{xx}
 -\frac{2(r+p)q}{cr^2}\psi_{xy}
 +\frac{a-q^2/r^2}{cr^2}\psi_{yy}\\
 &\hspace{20mm}
 -\frac{a-q^2/r^2}{cr}p
 -\frac{2(r+p)q^2}{cr^3}=0.\end{aligned}\] Consequently, the coefficient multiplying \(\psi_{yy}\) is \(1/c+O_3^{\rm corr}\), with \[O_3^{\rm corr}=
 \frac{(2c-x)x-(\gamma-1)(z+rp+p^2/2)
             -(\gamma+1)q^2/(2r^2)}{cr^2}.\] The printed expression instead has \(-p^2/2\) within the parenthesis multiplied by \(-(\gamma-1)\), and the wrong quadratic \(q\) contribution. Their difference is \[O_3^{\rm print}-O_3^{\rm corr}
   =\frac{(\gamma-1)p^2+\gamma q^2/r^2}{cr^2}.\] These are not interchangeable coefficients.\par
For an exact rational check, set \[\gamma=2,\quad c=2,\quad x=\tfrac12,\quad
 z=\tfrac1{16},\quad p=\tfrac14,\quad q=0,\qquad
 \psi_{xx}=\psi_{xy}=0,\quad\psi_{yy}=\tfrac38.\] Then \[a=\tfrac{113}{32},\quad |\mathrm D\varphi|^2=\tfrac{49}{16},
 \quad a-|\mathrm D\varphi|^2=\tfrac{15}{32}>0,
 \quad O_3^{\rm corr}=\tfrac{41}{144},\quad
 O_3^{\rm print}=\tfrac{43}{144}.\] The original transformed equation has residual zero, while (3.25) with the printed \(O_3\) has residual \(1/192\).\par
This jet can be realized by an actual local constant-state solution, not merely an arbitrary Hessian. At \(r_0=c-x\) and angle \(y_0\), take \[\psi(r,y)=-p\,r\cos(y-y_0)+z+pr_0.\] This is affine in Cartesian coordinates; \(\varphi=\varphi_5+\psi\) is a constant-state pseudo-potential. At the selected point its polar derivatives are exactly those above, including \(\psi_{yy}=r_0p\). Continuity gives a strictly elliptic neighborhood. This is a local equation counterexample, not a global free-boundary counterexample.\par
The defect persists arbitrarily near a sonic point and for any \(c>0\). For \(\gamma=2\), \(0<x<c\), take \(p=x/2\), \(z=x^2/4\), \(q=0\) and the same affine construction. Then \[a-|\mathrm D\varphi|^2=\frac{cx}{2}-\frac{x^2}{8}>0,
 \qquad \text{printed residual}=\frac{p^3}{cr}>0.\] Thus the discrepancy cannot be dismissed as a jet far outside a small sonic neighborhood.\par
Core impact. After cancelling removable factors in the formula for \(O_4\), the correct coefficient satisfies \[r(1/c+O_3^{\rm corr})=1+O_4,\qquad O_5=O_2/r.\] These identities cancel the first-order terms when converting the polar equation back to Cartesian coordinates. The printed coefficient leaves an extra multiple of \(\psi_x\): \[\left(\frac{(\gamma-1)p^2}{cr}
       +\frac{\gamma q^2}{cr^3}\right)\psi_x.\] The claim on p. 46 that these terms vanish is therefore false for the stated coefficients. A cutoff construction does not create the missing algebraic identity in the region where the cutoffs leave the original quantities unchanged.\par
Counterevidence and repair boundary. Replacing \(O_3\) by the derived expression restores the cancellation. For gradients of order \(x\), the difference is of quadratic order, so it does not change the leading coefficient responsible for sonic degeneracy. This supports the possibility of repairing the construction. However, the actual modified operator, its estimates, and the fixed-point recovery of the Euler equation must then be rechecked with the corrected coefficient. A repair formula alone does not validate those dependencies. The \(\gamma=1\) endpoint also requires cancelling the removable \((\gamma-1)\) factors in \(O_4\) before evaluation. The classification is Major because the as-printed core operator identification fails, not because a main theorem has been shown false.\par
\subsection*{F2. The moving-domain estimates omit essential dependence (Minor)}
Location. Definition 4.7, p. 41; Lemma 4.8(iv), p. 42, and the second estimate in Lemma 4.8(vi), p. 43.\par
The local \(C^{2,\alpha}(K)\) estimate for the differences of \(\varphi\circ F\) and \(\psi\circ F\) is bounded only by a local norm of \(u-\widetilde u\) on \(K\) and the parameter difference. However, \(F\) depends on the top boundary trace \(u(s,1)\), which is not determined by that interior norm.\par
Use the permitted normal-reflection extension at \(\theta=(0,0)\), \(\gamma=2\), \(v_2=-1\). Then \[\rho_0=\frac{5+\sqrt{17}}4,\qquad
 \eta_0=\frac{\sqrt{17}-1}4>0,\qquad
 \varphi_0=-\frac{|\boldsymbol\xi|^2}{2}-\eta_0.\] Let \(\widetilde u=0\) and \(u_\varepsilon(s,t)=\varepsilon a(s)b(t)\), with \(a\) smooth, compactly supported in the middle of \((-1,1)\) and \(a(0)=1\). Choose \(b=0\) near a compact interior set \(K\) containing \((0,t_0)\) and \(b(1)=1\). For sufficiently small \(\varepsilon\), the smallness requirement is satisfied; for each such smooth perturbation the required finite higher norm bound can be chosen. All side traces vanish.\par
The central fiber of the first flattening has \(G_1(0,\eta)=(0,\eta)\). Since \(\varphi_2-\varphi_0=\eta_0-\eta\), the perturbed shock height is \(g_\varepsilon(0)=\eta_0-\varepsilon\), and \[F_\varepsilon(0,t)=(0,t(\eta_0-\varepsilon)).\] Although \(u_\varepsilon-\widetilde u\) is identically zero near \(K\) and the parameter difference is zero, the potential difference at \((0,t_0)\) is \[(\varphi\circ F_\varepsilon-\widetilde\varphi\circ F_0)(0,t_0)
 =t_0^2\left(\eta_0\varepsilon-\frac{\varepsilon^2}{2}\right)>0
 \quad (0<\varepsilon<2\eta_0).\] The asserted right-hand side vanishes, contradicting the estimate.\par
There is also an omitted dependence in the \(C^{2,\alpha}\) bound for \(F\) in (iv): its constant is stated to depend only on \(\gamma,v_2,\theta_*\), although the hypothesis admits an arbitrary higher norm bound \(M\) in (4.19). With fixed small \(\varepsilon\) and a central cutoff \(\chi\), choose the top trace \[a_N(s)=\varepsilon N^{-1-\alpha}\chi(s)\cos Ns.\] The \(C^{1,\alpha}\) norms remain uniformly small, while \(|a_N''(0)|=\varepsilon N^{1-\alpha}\to\infty\). Near the central fiber the height is \(\eta_0-a_N(s)\), so the second derivative of the vertical component of \(F(s,t)\) equals \(-t a_N''(s)\). Every member has a finite admissible \(M_N\), but no common bound independent of \(M_N\) is possible.\par
Counterevidence and repair. Add control of the top trace, or a norm of \(F-\widetilde F\), to the local composition estimate; permit dependence on the relevant higher norm bound in (iv). The later argument uses globally controlled sequences, for which the missing trace information is available. This calculation refutes the stated estimates, not compactness of the corrected iteration map.\par
\subsection*{F3. The strict normal-sign statement needs positive density (Minor)}
Location. Definition 2.1 and the discussion following (2.3), p. 5; the admitted vacuum configuration, p. 7.\par
The weak-solution class permits vacuum, but the subsequent strict same-sign conclusion for the two normal derivatives does not. Let \(\gamma=2\), \(B=0\), with interface \(\xi=0\), and set \[\varphi_L=-\tfrac12(\xi^2+\eta^2)+\sqrt2\,\xi
 \quad(\xi<0),\qquad
 \varphi_R=-\tfrac12(\xi^2+\eta^2)
 \quad(\xi>0).\] The continuous piecewise smooth potential is locally Lipschitz. Bernoulli’s law gives \(\rho_L=0\) and \(\rho_R=1\). The mass equation holds in each half-plane and distributionally: the normal mass flux on either side of the interface is zero. The energy inequality allows equality at vacuum. Yet the normal derivatives are \(\sqrt2\) and \(0\), not two strictly same-sign nonzero numbers, and the gradient magnitude jumps.\par
Counterevidence and repair. If both traces of density are strictly positive, the Rankine–Hugoniot condition gives equal-sign normal derivatives. A zero common flux would force both to vanish; continuity of the tangential derivative would then rule out a shock with a gradient-magnitude jump. Thus adding positive densities supplies the claimed strict conclusion. The main nonvacuum parameter range excludes this example. The defect is in the general statement accompanying the weak formulation.\par
\subsection*{F4. The radial comparison uses the wrong gradient (Major)}
Location. Radial coordinates about the upstream sonic center, p. 24; identity (3.44), p. 32.\par
Let \(w=\varphi_2-\varphi\) and \(e=\mathrm D\varphi_2(P_0^1)/|\mathrm D\varphi_2(P_0^1)|\). Since the radial unit vector about the upstream center is \(-\mathrm D\varphi_2/|\mathrm D\varphi_2|\), the correct identity is \[\partial_e w=-\partial_r w+
 \left(e-\frac{\mathrm D\varphi_2}{|\mathrm D\varphi_2|}\right)\cdot\mathrm Dw.\] The last normalized vector in (3.44) is instead \(\mathrm D\varphi/|\mathrm D\varphi|\).\par
At a subsonic reference corner, write \(R=\xi_{P_0^1}\), \(u=u_5\), \(v=v_2\) and \(M=\sqrt{R^2+v^2}\), with \(0<u<R\) and \(v<0\). The reference gradients are \[\mathrm D\varphi_2=(-R,v),\quad \mathrm D\varphi=(-(R-u),0),\quad
 \mathrm Dw=(-u,v).\] The printed right-hand side minus the correct left-hand side is \[\frac{v^2}{M}\left(1-\frac{u}{M+R}\right)>0.\] The mismatch at the continuous endpoint also prevents the identity from holding on an adjacent open shock segment.\par
Counterevidence and repair. The nonzero discrepancy at the reference corner prevents the printed correction from being treated as a quantity tending to zero. Equation (3.44) is used immediately on p. 32 to obtain the positive directional derivative for hodograph coordinates. The resulting subsonic-corner estimates enter Theorem 2.2(i) on p. 60. Thus F4 is Major for an invalid derivation on the central regularity proof path.\par
Restoring the subscript \(2\) makes the error vector tend to zero; the available gradient and radial bounds then give the intended directional bound after shrinking the neighborhood. This supported repair does not validate the printed identity. No counterexample to hodograph monotonicity or the regularity theorem is asserted.\par
\subsection*{F5. Normal reflection exposes a strict inequality and a pullback sign error (Minor)}
Location. Lemma 4.17, Step 2, p. 52; the pullback introduced in the compactness proof on p. 58; normal-reflection calculation on p. 59.\par
At \((u,\theta)=(u^\sharp,\theta^\sharp)=(0,0)\), condition (4.40) holds with zero distance. The normal solution gives \[\widehat\phi=v_2\eta_0,\qquad
 \varphi_0^*=-|\boldsymbol\xi|^2/2+v_2\eta_0.\] Thus \(\widehat\phi-(\varphi_\theta^*+|\boldsymbol\xi|^2/2)=0\), contradicting the strict positive inequality asserted in Step 2 for every nearby member. Normal reflection belongs to the extended iteration set used in that statement.\par
On p. 58 the displayed definition of \(\widehat w\) uses \[(\widehat\phi+|\boldsymbol\xi|^2/2-\varphi_\theta^*)\circ G_1^{-1}.\] The sign before \(|\boldsymbol\xi|^2/2\) must be minus, consistently with (4.24) and (4.41). At normal reflection the printed expression is \(|\boldsymbol\xi|^2\) before composition, whereas the intended perturbation is identically zero.\par
Counterevidence and repair. Separate the exact normal solution from the nonzero-angle positivity argument, and use nonnegativity where the extended parameter range includes zero. Correct the pullback sign. The normal solution already satisfies (4.39) with zero perturbation norm, so this example does not refute that estimate or compactness.\par
\subsection*{F6. The endpoint differentiation and mean-value step are incorrect as written (Major)}
Location. Equation (5.3), p. 60; (5.5)–(5.6) and the intervening mean-value assertion, p. 61.\par
Put \(f=f_{5,\mathrm{sh}}\), \(g=f_{5,0}\) and \(Q=\overline\psi_{25}\). Differentiating \(Q(x,f(x))-Q(x,g(x))=\psi(x,f(x))\) twice gives \[\begin{aligned}
 Q_y(x,f)(f-g)''={}&
 \sum_{k=0}^2\binom2k
 \bigl[(g')^k\partial_x^{2-k}\partial_y^kQ(x,g)
       -(f')^k\partial_x^{2-k}\partial_y^kQ(x,f)\bigr]\\
 &+[Q_y(x,g)-Q_y(x,f)]g''
   +f''\psi_y(x,f)\\
 &+\sum_{k=0}^2\binom2k
     (f')^k\partial_x^{2-k}\partial_y^k\psi(x,f).\end{aligned}\] The printed terms use powers of \(f,g\) instead of their derivatives in the relevant chain-rule factors, show incorrect derivative indices in \(A_1\), and give the wrong sign to \(A_3\). A simple identity test is \[Q(x,y)=y,\qquad \psi(x,y)=x^2,\qquad f(x)=x^2,\qquad g(x)=0.\] It satisfies the starting identity and the matching conditions at zero. The left-hand side is \(2\), while the printed right-hand side is \(-2\); the corrected expression is \(2\). This is a test of the claimed differentiation identity, not an assertion that these functions solve the fluid problem.\par
The later implication from \(F(0)=0\) and \(F(x)/x\to0\) to positive points with \(F'(x_m)=0\) is also false: \(F(x)=x^2\) has \(F'(x)>0\) for every \(x>0\). The mean value theorem instead gives, for \(h_m\downarrow0\), some \(x_m\in(0,h_m)\) with \[F'(x_m)=F(h_m)/h_m\longrightarrow0.\] This weaker, correct conclusion suffices for the subsequent limiting argument. In (5.6), differentiation along the curve likewise requires the factor \(f'(x_m)-\omega/10\), not \(f(x_m)-\omega/10\).\par
Counterevidence and repair. On p. 60, the twice-differentiated identity (5.3) is used to deduce the matching second derivative and \(C^{2,\alpha}\) shock joining in Theorem 2.2(i). Its failure is Major for that core derivation. The identity test is not an Euler solution counterexample.\par
The separate exact-zero assertion on p. 61 has a narrower effect: the subsequent limit needs only \(F'(x_m)\to0\), already supplied by the mean value theorem under the given hypotheses. In (5.6), both printed and corrected factors are bounded and multiply \(\psi_{xy}\to0\). These overstatements do not supply another failure of the limiting conclusion.\par
The corrected twice-differentiated formula and the existing derivative, corner and vanishing bounds restore the specified endpoint argument. That repair does not lower the severity of the original derivation. The claimed optimal regularity has not been disproved.\par
\subsection*{F7. Two Bernoulli displays have incorrect algebra (Minor)}
Location. Lemma 3.3, Step 2, p. 19; the Bernoulli difference displayed on p. 24.\par
Set \(C=1+(\gamma-1)(B-\varphi)\). Then \[c^2=C-\frac{\gamma-1}{2}|\mathrm D\varphi|^2.\] Strict subsonicity implies \(|\mathrm D\varphi|^2<2C/(\gamma+1)\), but it does not imply the displayed equality \(c^2=2C/(\gamma+1)\) for \(\gamma>1\). For instance, \(\gamma=2\), \(C=1\), \(\mathrm D\varphi=0\) gives \(c^2=1\), not \(2/3\). The correct strict inequality follows by rearranging \(|\mathrm D\varphi|^2<c^2\) directly.\par
Since \(h(\rho_2)=h(1)=0\), subtraction of the two Bernoulli equations also gives \[|\mathrm D\varphi_2|^2-|\mathrm D\varphi|^2
 =2h(\rho)+2(\varphi-\varphi_2).\] The second term on the right of the p. 24 display is missing the factor \(2\).\par
Counterevidence and repair. The first correction gives the same upper bound needed for the Lipschitz estimate. For the second, the available Lipschitz control and equality \(\varphi=\varphi_2\) on the shock allow the neighborhood to be made small enough even with the corrected factor. Neither algebraic correction contradicts the intended positive-gap estimate.\par
\section*{Unresolved dependencies and rejected candidates}
Definition 4.11 of the extended iteration set on p. 45 retains (i)–(v) of the preceding definition and omits its ellipticity condition. Yet away from the sonic collars the matrix on p. 46 is \(c^2I-\mathrm D\varphi\otimes\mathrm D\varphi\), and later assertions require ellipticity for the full extended set. It is not immediate from the printed geometric, density and norm conditions that the omitted inequality follows. This is recorded as an unresolved proof dependency, not a confirmed error: no global member satisfying every retained condition and violating ellipticity was constructed. A proof of that implication or an explicit additional restriction would settle it.\par
Several apparent extraction-based problems were rejected after inspection of the actual page images. On pp. 21–22 the limiting parameter lies in the closed parameter region. In the weak formulation, test functions have support in the closed upper half-plane. The weighted spaces change their weight together with the Hölder exponent; an alleged compactness obstruction based on a fixed weight was therefore inapplicable. Endpoint traces that follow automatically from the weighted norm do not create the suggested nonopenness problem. None of these rejected candidates contributes to the assessment.\par
The proof uses external results for mixed boundary problems, degenerate Schauder estimates and shock convexity. The full original proofs of those cited results were not reacquired and reproduced in this review. This limitation must not be replaced by a declaration that every dependency is verified.\par
\section*{Supporting checks}
F1 invalidates the core equation identification, including its stated Cartesian cancellation. F4 invalidates the radial-comparison step used for subsonic-corner hodograph estimates. F6 invalidates the differentiated identity used for shock joining in Theorem 2.2(i); its separate exact-zero assertion is stronger than the limiting argument needs. These are the three supported Major findings. F2, F3, F5 and F7 remain Minor within their stated application and parameter boundaries.\par
The corrected radial and differentiation arguments restore the specified local steps. A corrected \(O_3\) is necessary, but closure of every subsequent estimate and the fixed-point construction still requires checking. Repairability does not reduce a core derivation error. No counterexample to Theorems 2.1–2.3 as existence, regularity or convexity assertions has been established.\par
The primary reviewer’s own SymPy calculation checks the coefficient difference, elliptic jet, normal mapping example, vacuum example and radial discrepancy; no author code was executed. Source rendering and actual image inspection are recorded separately. All superseded bounded-check drafts and corrected extraction candidates are retained.\par
A bounded publisher/arXiv and supplement/code search found no acquired paper-specific supplement or code artifact. Generic arXiv code widgets are not evidence of an associated program. Absence of an acquired code artifact is not a mathematical deficiency. The Oxford repository surfaced as an alternative source but its file was not acquired or compared. This review remains restricted to the frozen v1 file.\par
9 G.-Q. G. Chen, A. Cliffe, F. Huang, S. Liu and Q. Wang, Global solutions of the two-dimensional Riemann problem with four-shock interactions for the Euler equations for potential flow, arXiv:2305.15224v1 (2023), acquired 75-page PDF. arXiv record, https://arxiv.org/abs/2305.15224, inspected 27 October 2026; used only for version metadata and optional-material lookup. EMS Press article record, https://ems.press/journals/jems/articles/14298711, inspected 27 October 2026; journal DOI 10.4171/JEMS/1622. The mathematical journal text was not acquired.\par
\end{document}
